% 87-01 ①(87쪽) — 보기 그래프: $x=-2$ 에서 극소, $x=1$ 에서 기울기 0 인 아래로 볼록한 사차함수 개형. 그림 자체가 보기라 TikZ 로 다시 그림.
% 라벨은 원문 그대로 -2, O, 1, x, y, y=f(x) 와 점선 두 줄만. 어느 것이 답인지 드러나는 표시는 넣지 않는다.
\begin{tikzpicture}[baseline=(current bounding box.center), x=0.40cm, y=0.40cm, line width=0.5pt]
  \path (-3.6,-2.2) rectangle (3.45,3.6);
  \draw[->, >=stealth, line width=0.7pt] (-3.35,0) -- (2.75,0) node[below=1pt, font=\footnotesize] {$x$};
  \draw[->, >=stealth, line width=0.7pt] (0,-1.7) -- (0,3.45) node[left=1pt, font=\footnotesize] {$y$};
  \draw[densely dotted, line width=0.4pt] (-2,0) -- (-2,-1.2);
  \draw[densely dotted, line width=0.4pt] (1,0) -- (1,1.64);
  \draw[line width=0.6pt, smooth] plot coordinates {(-3.15,2.6) (-3.05,2.1) (-2.95,1.6) (-2.85,1.15) (-2.75,0.7) (-2.6,0.15) (-2.45,-0.35) (-2.3,-0.75) (-2.15,-1.05) (-2,-1.2) (-1.85,-1.15) (-1.7,-1.0) (-1.5,-0.7) (-1.3,-0.35) (-1.1,0) (-0.85,0.45) (-0.55,0.9) (-0.25,1.2) (0.1,1.45) (0.45,1.57) (0.8,1.62) (1,1.64) (1.3,1.7) (1.6,1.88) (1.9,2.2) (2.2,2.7)};
  \node[above=1pt, font=\footnotesize] at (-1.88,0) {$-2$};
  \node[below left=-2pt and -3pt, font=\footnotesize] at (0,0) {$\pt{O}$};
  \node[below=1pt, font=\footnotesize] at (1,0) {$1$};
  \node[anchor=west, font=\footnotesize] at (0.28,3.25) {$y=f(x)$};
\end{tikzpicture}
